\setcounter{ex}{0}

% --- CÂU 1 ---
\begin{ex}
	Gọi $M, m$ lần lượt là giá trị lớn nhất, giá trị nhỏ nhất của hàm số $y = \sin^2 x + 2\cos x$. Tính $M - m$.
	\shortans{4}
	\loigiai{
		Ta có:
		$$y = \sin^2 x + 2\cos x = 1 - \cos^2 x + 2\cos x = 2 - (\cos x - 1)^2.$$
		Với mọi $x \in \mathbb{R}$, vì $-1 \le \cos x \le 1$ nên $-2 \le \cos x - 1 \le 0 \implies 0 \le (\cos x - 1)^2 \le 4$.\\
		Suy ra:
		$$-2 \le 2 - (\cos x - 1)^2 \le 2 \implies -2 \le y \le 2.$$
		Do đó $M = 2$ (đạt được khi $\cos x = 1$) và $m = -2$ (đạt được khi $\cos x = -1$).\\
		Vậy $M - m = 2 - (-2) = 4$.
	}
\end{ex}

% --- CÂU 2 ---
\begin{ex}
	Cho $\sin\alpha = -\dfrac{4}{5}$ và $\pi < \alpha < \dfrac{3\pi}{2}$. Tính $A = 5\cos\alpha - 3\tan\alpha$.
	\shortans{-7}
	\loigiai{
		Vì $\pi < \alpha < \dfrac{3\pi}{2}$ nên $\cos\alpha < 0$. Ta có:
		$$\cos\alpha = -\sqrt{1 - \sin^2\alpha} = -\sqrt{1 - \left(-\dfrac{4}{5}\right)^2} = -\dfrac{3}{5}.$$
		Suy ra $\tan\alpha = \dfrac{\sin\alpha}{\cos\alpha} = \dfrac{-\frac{4}{5}}{-\frac{3}{5}} = \dfrac{4}{3}$.\\
		Thay vào biểu thức $A$, ta được:
		$$A = 5 \cdot \left(-\dfrac{3}{5}\right) - 3 \cdot \dfrac{4}{3} = -3 - 4 = -7.$$
	}
\end{ex}

% --- CÂU 3 ---
\begin{ex}
	Một hội trường có $15$ hàng ghế, hàng thứ nhất có $20$ ghế, mỗi hàng sau nhiều hơn hàng liền trước $3$ ghế. Hội trường có tất cả bao nhiêu ghế?
	\shortans{615}
	\loigiai{
		Số ghế ở mỗi hàng từ hàng thứ nhất đến hàng thứ $15$ lập thành một cấp số cộng với số hạng đầu $u_1 = 20$, công sai $d = 3$ và số các số hạng là $n = 15$.\\
		Tổng số ghế trong hội trường là:
		$$S_{15} = 15u_1 + \dfrac{15 \cdot 14}{2}d = 15 \cdot 20 + 105 \cdot 3 = 300 + 315 = 615\text{ (ghế)}.$$
	}
\end{ex}

% --- CÂU 4 ---
\begin{ex}
	Chiều cao sóng biển (mét) tại thời điểm $t$ (giờ) là $h(t) = 2 + 1{,}5\sin\left(\dfrac{\pi}{6}t\right)$. Tìm thời điểm đầu tiên $t > 0$ để chiều cao sóng bằng $2{,}75\text{ m}$.
	\shortans{1}
	\loigiai{
		Để chiều cao sóng bằng $2{,}75\text{ m}$ thì:
		$$h(t) = 2{,}75 \iff 2 + 1{,}5\sin\left(\dfrac{\pi}{6}t\right) = 2{,}75 \iff \sin\left(\dfrac{\pi}{6}t\right) = \dfrac{1}{2}$$
		$$\iff \hoac{\dfrac{\pi}{6}t &= \dfrac{\pi}{6} + k2\pi \\ \dfrac{\pi}{6}t &= \dfrac{5\pi}{6} + k2\pi} \iff \hoac{t &= 1 + 12k \\ t &= 5 + 12k}\ (k \in \mathbb{Z}).$$
		Thời điểm đầu tiên $t > 0$ ứng với $k = 0$ ở họ nghiệm thứ nhất là $t = 1$ (giờ).
	}
\end{ex}

% --- CÂU 5 ---
\begin{ex}
	Cho hình chóp $S.ABCD$ có đáy $ABCD$ là hình bình hành tâm $O$, $M$ là trung điểm của $SC$, $I = AM \cap SO$. Biết $\dfrac{SI}{SO} = \dfrac{a}{b}$ với $\dfrac{a}{b}$ là phân số tối giản. Tính $a + b$.
	\shortans{5}
	\loigiai{
		Vì $ABCD$ là hình bình hành tâm $O$ nên $O$ là trung điểm của $AC$.\\
		Xét tam giác $SAC$ có hai đường trung tuyến $SO$ và $AM$ cắt nhau tại $I$, suy ra $I$ là trọng tâm của tam giác $SAC$.\\
		Theo tính chất trọng tâm tam giác, ta có:
		$$\dfrac{SI}{SO} = \dfrac{2}{3}.$$
		Do $\dfrac{2}{3}$ là phân số tối giản nên $a = 2, b = 3 \implies a + b = 2 + 3 = 5$.
	}
\end{ex}

% --- CÂU 6 ---
\begin{ex}
	Một vật dao động điều hòa có li độ (cm) tại thời điểm $t$ (giây) là $x = 3\sin t + 4\cos t$. Tính biên độ dao động của vật (đơn vị: cm).
	\shortans{5}
	\loigiai{
		Ta có:
		$$x = 3\sin t + 4\cos t = \sqrt{3^2 + 4^2}\left(\dfrac{3}{5}\sin t + \dfrac{4}{5}\cos t\right) = 5\cos(t - \varphi),$$
		trong đó $\cos\varphi = \dfrac{4}{5}$ và $\sin\varphi = \dfrac{3}{5}$.\\
		Vậy biên độ dao động của vật là $A = \sqrt{3^2 + 4^2} = 5\text{ (cm)}$.
	}
\end{ex}
